% ==============================================================================
% PHỤ LỤC B: HƯỚNG DẪN CÀI ĐẶT VÀ VẬN HÀNH
% ==============================================================================

\chapter{HƯỚNG DẪN CÀI ĐẶT VÀ VẬN HÀNH}
\label{app:installation}

\section{Yêu cầu hệ thống tối thiểu}

Để triển khai và vận hành hệ thống PetCare AI, môi trường máy tính cần đáp ứng các yêu cầu sau:
\begin{itemize}[leftmargin=1.5cm]
    \item \textbf{Hệ điều hành:} Windows 10/11 64-bit (kích hoạt WSL 2), Ubuntu 22.04 LTS hoặc macOS 13+.
    \item \textbf{Phần mềm nền tảng:} Docker Desktop v2.20+ kết hợp Docker Compose, Node.js 20+ và npm 10+.
    \item \textbf{Bộ nhớ RAM:} Tối thiểu 8 GB RAM (khuyến nghị 16 GB để chạy đồng thời 4 containers và server phát triển web).
    \item \textbf{Dung lượng ổ đĩa:} Trống tối thiểu 10 GB SSD.
\end{itemize}

\section{Các bước cài đặt và khởi động hệ thống}

\subsection{Bước 1: Thiết lập biến môi trường}
Sao chép tệp mẫu \texttt{.env.example} và cấu hình khóa API Gemini của Google AI Studio:
\begin{lstlisting}[language=bash, caption={Tạo tệp cấu hình môi trường}]
# Sao chep tep cau hinh mau
cp .env.example .env

# Dien khoa GEMINI_API_KEY cua ban vao tep .env
GEMINI_API_KEY=AIzaSy...
AI_PROVIDER=gemini
GEMINI_MODEL=gemini-3.8-flash
\end{lstlisting}

\subsection{Bước 2: Khởi chạy cụm dịch vụ bằng Docker Compose}
Tại thư mục gốc dự án \texttt{petCare/}, thực thi câu lệnh điều phối duy nhất:
\begin{lstlisting}[language=bash, caption={Khởi chạy hệ thống 4 dịch vụ container}]
docker compose up -d
\end{lstlisting}
Quá trình trên tự động kích hoạt:
\begin{itemize}
    \item \texttt{petcare\_db}: PostgreSQL 15 (Cổng 5432)
    \item \texttt{petcare\_redis}: Redis 7 (Cổng 6379)
    \item \texttt{petcare\_backend}: FastAPI API Server (Cổng 8001)
    \item \texttt{petcare\_worker}: Celery Worker chạy ngầm lịch tiêm
\end{itemize}

\subsection{Bước 3: Khởi chạy Giao diện Frontend Web}
Mở một cửa sổ terminal mới và khởi chạy ứng dụng Next.js:
\begin{lstlisting}[language=bash, caption={Cài đặt gói và khởi chạy Next.js dev server}]
cd web
npm install
npm run dev
\end{lstlisting}

Sau khi máy chủ khởi động thành công, truy cập các địa chỉ tương ứng trên trình duyệt:
\begin{itemize}[leftmargin=1.5cm]
    \item \textbf{Giao diện chính:} \url{http://localhost:3000}
    \item \textbf{Siêu thị thú cưng:} \url{http://localhost:3000/shop}
    \item \textbf{Tài liệu API Swagger UI:} \url{http://localhost:8001/docs}
\end{itemize}
